\section{Kramers--Wannier duality squared}%
\label{sec:asymex_kw}

The unnormalized Kramers--Wannier duality kernel on a periodic chain of $L$
qubits is
\begin{align}
	\bra{s'}D\ket{s} & =\prod_{j=1}^{L}(-1)^{s'_j(s_j+s_{j+1})},
	\qquad s_{L+1}=s_1,
	\label{eq:asymex_kw_kernel}
\end{align}
the periodic operator of a matrix product operator of bond dimension two.
Squaring the duality gives
\begin{align}
	D^2 & =2^L(1+\eta)\,T,
	\label{eq:asymex_kw_fusion}
\end{align}
where $\eta=\prod_jX_j$ is the global spin flip and $T$ is the one-site
left translation. This is the lattice form of the duality-defect fusion
rule of the Ising model \cite{Aasen2016Topological}, in the normalization in
which the analogous Majorana-chain defect satisfies $D^2=(1+\eta)T$
\cite{Seiberg2024Majorana}. It is a fusion rule of a non-invertible
symmetry whose right-hand side contains a space-time symmetry, so it lies
outside every fusion-category algebra of matrix product operators. Both
$T$ and $\eta T$ are bond-dimension-two matrix product operators whose four
matrices are the four matrix units of $\MN{2}$, hence normal already at
word length one.

\begin{definition}[The duality kernel and the two translation tensors]
	\label{def:asymex_kw}
	\lean{KWExample.kwTensor, KWExample.shiftTensor, KWExample.flipShiftTensor,
		KWExample.kwSquare, KWExample.kwSquareTargets}
	\leanok
	\uses{def:mpo_tensor, def:mpo_to_mps, def:mpdo_operator_product}
	With the outgoing label $s'$ written first, the duality tensor is
	\begin{align}
		W^{00} & =\begin{pmatrix}1&1\\0&0\end{pmatrix},
		       &
		W^{01} & =\begin{pmatrix}0&0\\1&1\end{pmatrix},
		       &
		W^{10} & =\begin{pmatrix}1&-1\\0&0\end{pmatrix},
		       &
		W^{11} & =\begin{pmatrix}0&0\\-1&1\end{pmatrix}.
		\notag
	\end{align}
	The translation tensor is $T^{s's}=E_{s,s'}$ and the flipped
	translation tensor is $(\eta T)^{s's}=E_{s,1-s'}$, where $E_{pq}$ is
	the matrix unit with a single entry $1$ in row $p$ and column $q$. The
	source is the stacked product $B^{s's}=\sum_mW^{s'm}\otimes W^{ms}$ of
	bond dimension four, and the two targets are $2T$ and $2\eta T$.
\end{definition}

\begin{center}
	% Formula: B^{s's}=\sum_mW^{s'm}\otimes W^{ms}, the stacked product
	% of Definition~\ref{def:asymex_kw}.
	% Source: Definition~\ref{def:asymex_kw}.
	% Ink-to-index map: both rows are the duality tensor $W$; the shared
	% vertical pairleg at each site is the summed index $m$; the drawn
	% upper-row label $r$ is the per-site component of the outgoing
	% index $s'$.
	% Contraction set: the vertical pairlegs (index $m$) at every site
	% and the horizontal virtual bonds within each row; the physical
	% legs $(r,s)$, the per-site components of $(s',s)$, remain open.
	% Boundary signature: (virtual-west=0 | virtual-east=0 |
	% physical-up=3 | physical-down=3) on both sides.
	\tenkzkernel
	\begin{tenkz}[rows={op,op}, cols=4, boundary=periodic]
		\tn[skin=mpo, ports={90:physical:$r_1$, 270:physical}]{W} &
		\tn[skin=mpo, ports={90:physical:$r_2$, 270:physical}]{W} &
		\tn[skin=dots, ports={180:virtual, 0:virtual}]{} &
		\tn[skin=mpo, ports={90:physical:$r_N$, 270:physical}]{W} \\
		\tn[skin=mpo, ports={90:physical, 270:physical:$s_1$}]{W} &
		\tn[skin=mpo, ports={90:physical, 270:physical:$s_2$}]{W} &
		\tn[skin=dots, ports={180:virtual, 0:virtual}]{} &
		\tn[skin=mpo, ports={90:physical, 270:physical:$s_N$}]{W}
	\end{tenkz}
	\quad $=$ \quad
	\begin{tenkz}[rows={op}, cols=4, boundary=periodic]
		\tn[skin=mpo, ports={90:physical:$r_1$, 270:physical:$s_1$}]{B} &
		\tn[skin=mpo, ports={90:physical:$r_2$, 270:physical:$s_2$}]{B} &
		\tn[skin=dots, ports={180:virtual, 0:virtual}]{} &
		\tn[skin=mpo, ports={90:physical:$r_N$, 270:physical:$s_N$}]{B}
	\end{tenkz}
\end{center}

\begin{theorem}[Normality of the two translation tensors]
	\label{thm:asymex_kw_normal}
	\lean{KWExample.shiftMPS_isNormal, KWExample.flipShiftMPS_isNormal}
	\leanok
	\uses{def:asymex_kw, def:normal}
	The translation tensor and the flipped translation tensor are normal:
	the four matrices of each are the four matrix units of $\MN{2}$, which
	already span $\MN{2}$ at word length one.
\end{theorem}

\begin{proof}
	\leanok
	Every matrix unit occurs among the four matrices of each tensor, and
	the matrix units span $\MN{2}$.
\end{proof}

\begin{theorem}[Split compression of the squared duality]
	\label{thm:asymex_kw_compression}
	% Principal declaration: KWExample.kwSquare_isReduction.
	\lean{KWExample.kwSquare_compression, KWExample.kwSquare_isReduction,
		KWExample.kwSquare_left_mul_right_of_ne, KWExample.kwSquare_dim_eq,
		KWExample.kwSquare_remainder_eq_zero}
	\leanok
	\uses{def:asymex_kw, thm:asym_multiblock}
	The source of Definition~\ref{def:asymex_kw} carries the block
	triangular form of Theorem~\ref{thm:asym_multiblock} for the two slots
	$2T$ and $2\eta T$ with no zero slots, $4=2+2$. The compression pairs
	are mutually biorthogonal and satisfy $W_sB^wV_s=C_s^w$ for every word,
	and the remainder $R^{s's}=B^{s's}-\sum_sV_sC_s^{s's}W_s$ vanishes
	identically: the extension splits.
\end{theorem}

\begin{proof}
	\leanok
	\uses{thm:asymex_explicit_gauge}
	The change of bond coordinates is
	\begin{align}
		G & =\frac12\begin{pmatrix}
			            1 & 0 & 1  & 0  \\
			            0 & 1 & 0  & -1 \\
			            1 & 0 & -1 & 0  \\
			            0 & 1 & 0  & 1
		            \end{pmatrix},
		\label{eq:asymex_kw_gauge}
	\end{align}
	whose inverse has integer entries. Conjugating each of the four letters
	by $G$ produces the block diagonal matrix with the blocks $2T$ and
	$2\eta T$, which after clearing the denominator is a finite verification
	over integer matrices. Block diagonality gives both the triangularity
	condition and the vanishing of the remainder, and the dimension count is
	clause (vii) of Theorem~\ref{thm:asym_multiblock}.
\end{proof}

\begin{theorem}[Periodic data of the squared duality]
	\label{thm:asymex_kw_trace}
	\lean{KWExample.kwSquare_trace_evalWord}
	\leanok
	\uses{def:asymex_kw}
	For every nonempty word $w$ over the alphabet of index pairs,
	\begin{align}
		\tr(B^w) & =2^{|w|}\bigl(\tr(T^w)+\tr((\eta T)^w)\bigr),
		\label{eq:asymex_kw_trace}
	\end{align}
	which is (\ref{eq:asymex_kw_fusion}) at the level of periodic
	coefficients.
\end{theorem}

\begin{proof}
	\leanok
	\uses{thm:asymex_kw_compression, prop:asym_compression_trace}
	Proposition~\ref{prop:asym_compression_trace} applied to the datum of
	Theorem~\ref{thm:asymex_kw_compression} gives
	$\tr(B^w)=\tr((2T)^w)+\tr((2\eta T)^w)$, and pulling the weight $2$ out
	of each word evaluation contributes the factor $2^{|w|}$.
\end{proof}

\begin{theorem}[The sitewise intertwiners of the split case]
	\label{thm:asymex_kw_intertwiners}
	% Principal declaration: KWExample.kwSquare_mul_right0_eq.
	\lean{KWExample.kwSquareLeft0, KWExample.kwSquareLeft1,
		KWExample.kwSquareRight0, KWExample.kwSquareRight1,
		KWExample.kwSquare_mul_right0_eq, KWExample.kwSquare_mul_right1_eq,
		KWExample.left0_mul_kwSquare_eq, KWExample.left1_mul_kwSquare_eq,
		KWExample.kwSquareLeft0_mul_right0, KWExample.kwSquareLeft1_mul_right1}
	\leanok
	\uses{def:asymex_kw, thm:asymex_kw_compression}
	Write
	\begin{align}
		V_1 & =\begin{pmatrix}
			       1 & 0  \\
			       0 & 1  \\
			       1 & 0  \\
			       0 & -1
		       \end{pmatrix},
		    & V_2
		    & =\begin{pmatrix}
			       1  & 0 \\
			       0  & 1 \\
			       -1 & 0 \\
			       0  & 1
		       \end{pmatrix},
		    & W_1
		    & =\frac12\begin{pmatrix}
			              1 & 0 & 1 & 0  \\
			              0 & 1 & 0 & -1
		              \end{pmatrix},
		    & W_2
		    & =\frac12\begin{pmatrix}
			              1 & 0 & -1 & 0 \\
			              0 & 1 & 0  & 1
		              \end{pmatrix}.
		\notag
	\end{align}
	Then $W_1V_1=W_2V_2=\id$, and for every index pair,
	$B^{s's}V_1=V_1(2T)^{s's}$, $B^{s's}V_2=V_2(2\eta T)^{s's}$,
	$W_1B^{s's}=(2T)^{s's}W_1$ and $W_2B^{s's}=(2\eta T)^{s's}W_2$.
\end{theorem}

\begin{proof}
	\leanok
	\uses{thm:asymex_explicit_gauge}
	These four matrices are the rows of $G$ and the columns of $G^{-1}$ at
	the coordinates of the two slots, by
	Theorem~\ref{thm:asymex_explicit_gauge}. Each of the four intertwining
	identities is an equality between products of integer matrices after
	clearing the denominator, and the two normalizations are clause (iv) of
	Theorem~\ref{thm:asym_multiblock}.
\end{proof}

\begin{center}
	% Formula: the four sitewise identities of
	% Theorem~\ref{thm:asymex_kw_intertwiners}: B^{s's}V_1=V_1(2T)^{s's},
	% B^{s's}V_2=V_2(2\eta T)^{s's}, W_1B^{s's}=(2T)^{s's}W_1, and
	% W_2B^{s's}=(2\eta T)^{s's}W_2.
	% Source: Theorem~\ref{thm:asymex_kw_intertwiners}.
	% Ink-to-index map: the ring is the compression map $V_1$, $W_1$,
	% $V_2$ or $W_2$; the square node is the source $B$ or its target
	% $2T$ or $2\eta T$.
	% Contraction set: the single virtual bond joining the ring to the
	% square node in each pair; the physical pair leg remains open.
	% Boundary signature: (virtual-west=1 | virtual-east=1 |
	% physical-up=1 | physical-down=0) for each pictured equality.
	\tenkzkernel
	\begin{tenkz}[rows={wire}, cols=2, west=open, east=open]
		\tn[label pos=s, ports={90:physical:$s's$}]{B} & \tn[skin=ring]{V_1}
	\end{tenkz}
	\quad $=$ \quad
	\begin{tenkz}[rows={wire}, cols=2, west=open, east=open]
		\tn[skin=ring]{V_1} & \tn[label pos=s, ports={90:physical:$s's$}]{2T}
	\end{tenkz}
	\\[2ex]
	\begin{tenkz}[rows={wire}, cols=2, west=open, east=open]
		\tn[label pos=s, ports={90:physical:$s's$}]{B} & \tn[skin=ring]{V_2}
	\end{tenkz}
	\quad $=$ \quad
	\begin{tenkz}[rows={wire}, cols=2, west=open, east=open]
		\tn[skin=ring]{V_2} &
		\tn[label pos=s, ports={90:physical:$s's$}]{2\eta T}
	\end{tenkz}
	\\[2ex]
	\begin{tenkz}[rows={wire}, cols=2, west=open, east=open]
		\tn[skin=ring]{W_1} & \tn[label pos=s, ports={90:physical:$s's$}]{B}
	\end{tenkz}
	\quad $=$ \quad
	\begin{tenkz}[rows={wire}, cols=2, west=open, east=open]
		\tn[label pos=s, ports={90:physical:$s's$}]{2T} & \tn[skin=ring]{W_1}
	\end{tenkz}
	\\[2ex]
	\begin{tenkz}[rows={wire}, cols=2, west=open, east=open]
		\tn[skin=ring]{W_2} & \tn[label pos=s, ports={90:physical:$s's$}]{B}
	\end{tenkz}
	\quad $=$ \quad
	\begin{tenkz}[rows={wire}, cols=2, west=open, east=open]
		\tn[label pos=s, ports={90:physical:$s's$}]{2\eta T} &
		\tn[skin=ring]{W_2}
	\end{tenkz}
\end{center}

\section{Kramers--Wannier duality acting on ground states}%
\label{sec:asymex_kw_action}

The same duality kernel may be applied to a state rather than to itself.
Both parts below read as symmetry breaking of the non-invertible
Kramers--Wannier symmetry: the duality maps the paramagnetic product state
to the pair of ferromagnetic sectors, and it maps both ferromagnetic
sectors to the paramagnet, merged with multiplicity two. When the acted
state is written as a multiple of a normalized state, a factor $2^{L/2}$
appears; that factor is a norm of the unnormalized periodic vector and not
a weight of the target, and the weighted tensor seen by the theorem is the
same under either reading.

\begin{definition}[The two action tensors]
	\label{def:asymex_kw_action}
	\lean{KWExample.plusTensor, KWExample.ghz0, KWExample.ghz1,
		KWExample.kwPlus, KWExample.plusTarget, KWExample.ghz,
		KWExample.kwGHZ, KWExample.kwGHZTarget}
	\leanok
	\uses{def:asymex_kw, def:mps_tensor}
	The product-state tensor is the bond-dimension-one tensor with both
	matrices equal to $(1)$, generating $\sum_s\ket{s}$ at every site; the
	two ferromagnetic sectors are the bond-dimension-one tensors
	$C_0=\bigl((1),(0)\bigr)$ and $C_1=\bigl((0),(1)\bigr)$, and their direct
	sum is the Greenberger--Horne--Zeilinger tensor $\diag(1,0)$,
	$\diag(0,1)$.
	The two action tensors are
	$(W\cdot A)^{s'}=\sum_sW^{s's}\otimes A^s$ applied to the product-state
	tensor and to the Greenberger--Horne--Zeilinger tensor; explicitly
	\begin{align}
		(W\cdot A)^0
		 & =\begin{pmatrix}1&1\\1&1\end{pmatrix},
		 &
		(W\cdot A)^1
		 & =\begin{pmatrix}1&-1\\-1&1\end{pmatrix}
		\notag
	\end{align}
	in the first case, with targets the two sectors carried with weight $2$,
	and
	\begin{align}
		(W\cdot A)^0
		 & =\begin{pmatrix}
			    1 & 0 & 1 & 0 \\
			    0 & 0 & 0 & 0 \\
			    0 & 0 & 0 & 0 \\
			    0 & 1 & 0 & 1
		    \end{pmatrix},
		 &
		(W\cdot A)^1
		 & =\begin{pmatrix}
			    1 & 0  & -1 & 0 \\
			    0 & 0  & 0  & 0 \\
			    0 & 0  & 0  & 0 \\
			    0 & -1 & 0  & 1
		    \end{pmatrix}
		\notag
	\end{align}
	in the second, with two unweighted copies of the product-state tensor as
	targets.
\end{definition}

\begin{center}
	% Formula: (W\cdot A)^{s'}=\sum_sW^{s's}\otimes A^s, of
	% Definition~\ref{def:asymex_kw_action}, evaluated at the
	% product-state tensor $A_+$ (left) and at the Greenberger--Horne--
	% Zeilinger tensor $A_{\mathrm{GHZ}}$ (right).
	% Source: Definition~\ref{def:asymex_kw_action}.
	% Ink-to-index map: the upper row is the duality tensor $W$, whose
	% drawn physical legs $r$ are the per-site components of the
	% outgoing index $s'$; the lower row is the acted-on state tensor;
	% the shared vertical pairleg at each site is the summed index $s$.
	% Contraction set: the vertical pairlegs (index $s$) at every site
	% and the horizontal virtual bonds within each row, both closed by
	% the periodic trace; only the upper physical legs of $W$ remain
	% open.
	% Boundary signature: (virtual-west=0 | virtual-east=0 |
	% physical-up=3 | physical-down=0) for both pictures.
	\tenkzkernel
	\begin{tenkz}[rows={op,ket}, cols=4, boundary=periodic]
		\tn[skin=mpo, ports={90:physical:$r_1$}]{W} &
		\tn[skin=mpo, ports={90:physical:$r_2$}]{W} &
		\tn[skin=dots, ports={180:virtual, 0:virtual}]{} &
		\tn[skin=mpo, ports={90:physical:$r_N$}]{W} \\
		\tn[label pos=s, ports={180:virtual, 0:virtual}]{A_+} &
		\tn[label pos=s, ports={180:virtual, 0:virtual}]{A_+} &
		\tn[skin=dots, ports={180:virtual, 0:virtual}]{} &
		\tn[label pos=s, ports={180:virtual, 0:virtual}]{A_+}
	\end{tenkz}
	\qquad
	\begin{tenkz}[rows={op,ket}, cols=4, boundary=periodic]
		\tn[skin=mpo, ports={90:physical:$r_1$}]{W} &
		\tn[skin=mpo, ports={90:physical:$r_2$}]{W} &
		\tn[skin=dots, ports={180:virtual, 0:virtual}]{} &
		\tn[skin=mpo, ports={90:physical:$r_N$}]{W} \\
		\tn[label pos=s, ports={180:virtual, 0:virtual}]{A_{\mathrm{GHZ}}} &
		\tn[label pos=s, ports={180:virtual, 0:virtual}]{A_{\mathrm{GHZ}}} &
		\tn[skin=dots, ports={180:virtual, 0:virtual}]{} &
		\tn[label pos=s, ports={180:virtual, 0:virtual}]{A_{\mathrm{GHZ}}}
	\end{tenkz}
\end{center}

\begin{theorem}[The duality on the product state splits]
	\label{thm:asymex_kw_plus}
	% Principal declaration: KWExample.plus_isReduction.
	\lean{KWExample.plusCompression, KWExample.plus_isReduction,
		KWExample.plus_dim_eq, KWExample.plus_remainder_eq_zero,
		KWExample.kwPlus_trace_evalWord_eq_sum}
	\leanok
	\uses{def:asymex_kw_action, thm:asym_multiblock}
	The first action tensor of Definition~\ref{def:asymex_kw_action} carries
	the block triangular form of Theorem~\ref{thm:asym_multiblock} for the
	two sectors carried with weight $2$, with no zero slots and $2=1+1$;
	the compression pair of each sector satisfies $W_sV_s=\id$ and
	$W_sB^wV_s=C_s^w$ for every word; the remainder vanishes identically;
	and for every nonempty word $w$,
	$\tr(B^w)=\tr((2C_0)^w)+\tr((2C_1)^w)$.
\end{theorem}

\begin{proof}
	\leanok
	\uses{thm:asymex_explicit_gauge, prop:asym_compression_trace}
	Conjugating the two letters by the matrix
	$\begin{pmatrix}1&1\\1&-1\end{pmatrix}$, whose inverse is half of
	itself, produces $\diag(2,0)$ and $\diag(0,2)$, a finite verification
	over integer matrices. Diagonality gives the triangularity condition,
	the identification of the matched blocks and the vanishing of the
	remainder; the trace identity is
	Proposition~\ref{prop:asym_compression_trace}.
\end{proof}

\begin{theorem}[The duality on the ferromagnetic state does not split]
	\label{thm:asymex_kw_ghz}
	% Principal declaration: KWExample.kwGHZ_isReduction.
	\lean{KWExample.kwGHZCompression, KWExample.kwGHZ_isReduction,
		KWExample.kwGHZ_dim_eq, KWExample.kwGHZ_evalWord_remainder_eq_zero,
		KWExample.kwGHZRemainder_eq, KWExample.kwGHZ_remainder_mul_remainder,
		KWExample.kwGHZ_trace_evalWord_eq_two}
	\leanok
	\uses{def:asymex_kw_action, thm:asym_multiblock}
	The second action tensor of Definition~\ref{def:asymex_kw_action}
	carries the block triangular form of Theorem~\ref{thm:asym_multiblock}
	for two unweighted copies of the product-state tensor with $z=2$ zero
	slots and $4=1+1+2$; the compression pair of each copy satisfies
	$W_sV_s=\id$ and $W_sB^wV_s=C_s^w$ for every word;
	$\tr(B^w)=2$ for every nonempty word $w$; a product of four remainder
	matrices vanishes; and in fact
	\begin{align}
		R^0 & =\begin{pmatrix}
			       0 & 0 & 1 & 0 \\
			       0 & 0 & 0 & 0 \\
			       0 & 0 & 0 & 0 \\
			       0 & 1 & 0 & 0
		       \end{pmatrix},
		    & R^1
		    & =\begin{pmatrix}
			       0 & 0  & -1 & 0 \\
			       0 & 0  & 0  & 0 \\
			       0 & 0  & 0  & 0 \\
			       0 & -1 & 0  & 0
		       \end{pmatrix},
		\notag
	\end{align}
	and $R^iR^j=0$ for all letters $i$ and $j$.
\end{theorem}

\begin{proof}
	\leanok
	\uses{thm:asymex_explicit_gauge, prop:asym_compression_trace}
	The gauge is the coordinate permutation placing the two copies at the
	first and fourth bond coordinates and the two zero slots at the second
	and third; in these coordinates the two letters are block upper
	triangular with the scalar $1$ on both matched diagonal positions, a
	finite verification over integer matrices. Every word evaluation of the
	product-state tensor is $1$, so the trace identity of
	Proposition~\ref{prop:asym_compression_trace} reads $\tr(B^w)=2$. The
	remainder is the conjugated letter off the diagonal of the block
	decomposition, which gives the two displayed matrices, and their
	pairwise products vanish by direct multiplication; the generic bound of
	clause (vi) of Theorem~\ref{thm:asym_multiblock} is $|S|+z=4$.
\end{proof}

\begin{theorem}[One-sided sitewise intertwiners for the merged target]
	\label{thm:asymex_kw_ghz_intertwiner}
	\lean{KWExample.kwGHZ_left_intertwiner_eq_zero,
		KWExample.kwGHZ_right_intertwiner,
		KWExample.kwGHZ_right_intertwiner_ne_zero}
	\leanok
	\uses{def:asymex_kw_action}
	If $u\in\C^4$ satisfies $uB^i=u$ for both letters $i$, then $u=0$: no
	nonzero sitewise left intertwiner exists for the product-state target.
	A nonzero sitewise right intertwiner does exist, namely the first
	coordinate vector, which satisfies $B^ie_0=e_0$ for both letters.
\end{theorem}

\begin{proof}
	\leanok
	Both matrices of the product-state target are the scalar $1$, so the
	two conditions read $uB^i=u$ and $B^iv=v$. Comparing the second and
	third coordinates of the two equations $uB^0=u$ and $uB^1=u$ gives in
	turn $u_3=u_1=u_0=u_2=0$. For the right-hand side, the first column of
	each of the two letters is the first coordinate vector.
\end{proof}

\section{Fibonacci fusion}%
\label{sec:asymex_fibonacci}

The simplest non-invertible fusion of matrix product operators is the
Fibonacci string-net algebra of \cite[Appendix~D.1]{Bultinck2017Anyons},
whose two operator families obey
\begin{align}
	O_\tau O_\tau & =O_1+O_\tau,
	\label{eq:asymex_fib_fusion}
\end{align}
the realization of $\tau\otimes\tau=1\oplus\tau$. The unit of the algebra
is not the identity operator: $O_1$ is the projector onto the admissible
configurations of the chain, those with no two neighbouring trivial labels,
and it is a matrix product operator of bond dimension two. The $\tau$
family has bond dimension three, so the stacked tensor of the square has
bond dimension nine: two directions carry the trivial block, three carry
the $\tau$ block, and the remaining four are annihilated by every letter,
being the inadmissible configurations killed by the vertex weights.

The entries of these tensors are not rational, so the exact arithmetic
takes place in a quartic ring.

\begin{definition}[The ring of golden integers]
	\label{def:asymex_golden_ring}
	\lean{GoldenInt, GoldenInt.sigma, goldenPoly, goldenSigmaReal,
		goldenToComplex, MPSTensor.complexOfGolden, MPSTensor.evalWordGolden,
		MPSTensor.mulGoldenTensor}
	\leanok
	Let $\sigma=\varphi^{-1/2}$ be the inverse square root of the golden
	ratio, the positive real root of $X^4+X^2-1$, so that
	$\sigma=\sqrt{(\sqrt5-1)/2}$. The ring of \emph{golden integers} is
	$\Z[\sigma]$, presented by the four integer coordinates of
	$c_0+c_1\sigma+c_2\sigma^2+c_3\sigma^3$ with the multiplication obtained
	from $\sigma^4=1-\sigma^2$. It embeds into $\C$ through the quotient
	$\mathbb{Q}[X]/(X^4+X^2-1)$ evaluated at that real root; a matrix over
	$\Z[\sigma]$ becomes a complex matrix by applying that embedding
	entrywise, and word evaluation and the bond-space product of tensors
	over $\Z[\sigma]$ are defined by the same formulas as over $\C$.
\end{definition}

\begin{theorem}[Exact arithmetic in the ring of golden integers]
	\label{thm:asymex_golden_ring}
	% Principal declaration: goldenMulRule.
	\lean{GoldenInt.sigma_quartic, goldenMulRule, goldenSigmaReal_quartic,
		MPSTensor.complexOfGolden_mul, MPSTensor.evalWord_complexOfGolden,
		MPSTensor.mulTensor_complexOfGolden}
	\leanok
	\uses{def:asymex_golden_ring}
	The generator satisfies $\sigma^4+\sigma^2-1=0$ in $\Z[\sigma]$, and so
	does the real number $\sqrt{(\sqrt5-1)/2}$. In any commutative ring
	carrying a root $s$ of $X^4+X^2-1$, the product of two integer
	combinations of $1,s,s^2,s^3$ is the combination prescribed by the
	multiplication of $\Z[\sigma]$. The entrywise embedding of matrices
	over $\Z[\sigma]$ into the complex matrices is multiplicative, and it
	commutes with word evaluation and with the bond-space product of
	tensors.
\end{theorem}

\begin{proof}
	\leanok
	Squaring $\sigma^2=(\sqrt5-1)/2$ gives the quartic for the real root.
	The multiplication rule is the reduction of a product of two
	polynomials of degree at most three modulo $X^4+X^2-1$, which is a
	polynomial identity in the eight coefficients. Multiplicativity of the
	entrywise embedding is the fact that it is induced by a ring
	homomorphism, and the two compatibility statements follow by induction
	on the word and by expanding the bond-space product entrywise.
\end{proof}

\begin{definition}[The two Fibonacci blocks and the stacked square]
	\label{def:asymex_fib}
	\lean{FibonacciCompression.fibTau, FibonacciCompression.fibOne,
		FibonacciCompression.fibTauMPS, FibonacciCompression.fibOneMPS,
		FibonacciCompression.fibStack, FibonacciCompression.fibTargets}
	\leanok
	\uses{def:asymex_golden_ring, def:mpo_tensor, def:mpdo_operator_product}
	Label the two objects of the category by $0$ for the trivial label and
	$1$ for $\tau$, and read a letter as the pair $(x',x)$ of the outgoing
	and incoming labels. Write
	\begin{align}
		F & =\begin{pmatrix}
			     0 & \varphi^{-1}   & \varphi^{-1/2} \\
			     1 & 0              & 1              \\
			     1 & \varphi^{-1/2} & -\varphi^{-1}
		     \end{pmatrix}
		\label{eq:asymex_fib_vertex}
	\end{align}
	for the reduced vertex weight. The $\tau$ tensor $B_\tau$ of bond
	dimension three has, at the letters $(0,1)$, $(1,0)$ and $(1,1)$, the
	matrix whose first, second and third row respectively is the
	corresponding row of $F$ and whose other rows vanish; at the letter
	$(0,0)$ it vanishes. The trivial tensor $B_1$ of bond dimension two
	carries the admissibility matrix
	$\begin{pmatrix}0&1\\1&1\end{pmatrix}$ on its two diagonal letters and
	vanishes on the other two. The source is the stacked product
	$B^{x'x}=\sum_m(B_\tau)^{x'm}\otimes(B_\tau)^{mx}$ of bond dimension
	nine, and the two targets are $B_1$ and $B_\tau$, each with
	multiplicity one and weight one.
\end{definition}

\begin{center}
	% Formula: B^{x'x}=\sum_m(B_\tau)^{x'm}\otimes(B_\tau)^{mx}, the
	% stacked product of Definition~\ref{def:asymex_fib}, realizing
	% $\tau\otimes\tau$; Theorem~\ref{thm:asymex_fib_fusion} identifies
	% its periodic trace with that of $B_1\oplus B_\tau$, the
	% realization $1\oplus\tau$.
	% Source: Definition~\ref{def:asymex_fib}.
	% Ink-to-index map: both rows are the $\tau$ tensor $B_\tau$; the
	% shared vertical pairleg at each site is the summed intermediate
	% label $m$; the drawn upper-row label $y$ is the per-site
	% component of the outgoing index $x'$.
	% Contraction set: the vertical pairlegs (index $m$) at every site
	% and the horizontal virtual bonds within each row; the physical
	% legs $(y,x)$, the per-site components of $(x',x)$, remain open.
	% Boundary signature: (virtual-west=0 | virtual-east=0 |
	% physical-up=3 | physical-down=3) on both sides.
	\tenkzkernel
	\begin{tenkz}[rows={op,op}, cols=4, boundary=periodic]
		\tn[skin=mpo, ports={90:physical:$y_1$, 270:physical}]{B_\tau} &
		\tn[skin=mpo, ports={90:physical:$y_2$, 270:physical}]{B_\tau} &
		\tn[skin=dots, ports={180:virtual, 0:virtual}]{} &
		\tn[skin=mpo, ports={90:physical:$y_N$, 270:physical}]{B_\tau} \\
		\tn[skin=mpo, ports={90:physical, 270:physical:$x_1$}]{B_\tau} &
		\tn[skin=mpo, ports={90:physical, 270:physical:$x_2$}]{B_\tau} &
		\tn[skin=dots, ports={180:virtual, 0:virtual}]{} &
		\tn[skin=mpo, ports={90:physical, 270:physical:$x_N$}]{B_\tau}
	\end{tenkz}
	\quad $=$ \quad
	\begin{tenkz}[rows={op}, cols=4, boundary=periodic]
		\tn[skin=mpo, ports={90:physical:$y_1$, 270:physical:$x_1$}]{B} &
		\tn[skin=mpo, ports={90:physical:$y_2$, 270:physical:$x_2$}]{B} &
		\tn[skin=dots, ports={180:virtual, 0:virtual}]{} &
		\tn[skin=mpo, ports={90:physical:$y_N$, 270:physical:$x_N$}]{B}
	\end{tenkz}
\end{center}

\begin{theorem}[The one-site matrices of the $\tau$ block]
	\label{thm:asymex_fib_entries}
	\lean{FibonacciCompression.fibTauMPS_apply_one,
		FibonacciCompression.fibTauMPS_apply_three}
	\leanok
	\uses{def:asymex_fib}
	The matrix of $B_\tau$ at the letter $(0,1)$ has the first row of
	(\ref{eq:asymex_fib_vertex}) as its first row and zero elsewhere, and
	its matrix at the letter $(1,1)$ has the third row of
	(\ref{eq:asymex_fib_vertex}) as its third row and zero elsewhere, with
	$\varphi^{-1/2}=\sqrt{(\sqrt5-1)/2}$.
\end{theorem}

\begin{proof}
	\leanok
	\uses{thm:asymex_golden_ring}
	Both are the entrywise embedding of the corresponding matrix over
	$\Z[\sigma]$, evaluated coordinate by coordinate.
\end{proof}

\begin{theorem}[Normality of the two Fibonacci blocks]
	\label{thm:asymex_fib_normal}
	\lean{FibonacciCompression.fibOneMPS_isNormal,
		FibonacciCompression.fibTauMPS_isNormal}
	\leanok
	\uses{def:asymex_fib, def:normal}
	Both blocks are normal, their length-three words spanning the full
	matrix algebra of their bond space.
\end{theorem}

\begin{proof}
	\leanok
	\uses{thm:asymex_golden_ring}
	For the $\tau$ block the word $(i+1,3,k+1)$, with $i$ and $k$ in
	$\{0,1,2\}$, evaluates to a multiple of the matrix built from the $i$-th
	coordinate vector and the $k$-th row of (\ref{eq:asymex_fib_vertex});
	those rows are a basis, so each of the nine matrix units is an explicit
	combination over $\Z[\sigma]$ of the nine such words. For the trivial block four length-three words suffice. Both
	families of coefficients are verified in $\Z[\sigma]$ and transported to
	the complex matrices by Theorem~\ref{thm:asymex_golden_ring}.
\end{proof}

\begin{theorem}[Split compression of the Fibonacci square]
	\label{thm:asymex_fib_compression}
	% Principal declaration: FibonacciCompression.fibonacci_isReduction.
	\lean{FibonacciCompression.fibonacci_compression,
		FibonacciCompression.fibonacci_isReduction,
		FibonacciCompression.fibonacci_left_mul_right_of_ne,
		FibonacciCompression.fibonacci_dim_eq,
		FibonacciCompression.fibonacci_evalWord_remainder_eq_zero,
		FibonacciCompression.fibonacci_remainder_eq_zero}
	\leanok
	\uses{def:asymex_fib, thm:asym_multiblock}
	The source of Definition~\ref{def:asymex_fib} carries the block
	triangular form of Theorem~\ref{thm:asym_multiblock} for the trivial
	block and the $\tau$ block with $z=4$ zero slots and $9=2+3+4$. The
	compression pairs are mutually biorthogonal and satisfy
	$W_sB^wV_s=C_s^w$ for every word; a product of six remainder matrices
	vanishes; and in fact the remainder vanishes identically, so the
	extension splits.
\end{theorem}

\begin{proof}
	\leanok
	\uses{thm:asymex_explicit_gauge, thm:asymex_golden_ring}
	The change of bond coordinates has as its columns the two fusion
	tensors of the blocks followed by a basis of the four-dimensional
	subspace annihilated by every letter. In these coordinates every letter
	of the stacked tensor is block diagonal with the trivial block, the
	$\tau$ block and four one-dimensional zero blocks on the diagonal, which
	is a finite verification over matrices with entries in $\Z[\sigma]$.
	Block diagonality gives the triangularity condition, the identification
	of the diagonal blocks and the vanishing of the remainder; the generic
	nilpotency bound of clause (vi) of Theorem~\ref{thm:asym_multiblock} is
	$|S|+z=6$, and the dimension count is clause (vii).
\end{proof}

\begin{theorem}[The fusion rule]
	\label{thm:asymex_fib_fusion}
	% Principal declaration: FibonacciCompression.fibonacci_fusion_rule.
	\lean{FibonacciCompression.fibStack_trace_evalWord,
		FibonacciCompression.fibonacci_fusion_rule,
		FibonacciCompression.fibStack_mpv}
	\leanok
	\uses{def:asymex_fib, def:mpo_family}
	For every nonempty word $w$ over the alphabet of label pairs,
	$\tr(B^w)=\tr(B_1^w)+\tr(B_\tau^w)$; equivalently, at every positive
	system size $L$ the periodic operators satisfy
	(\ref{eq:asymex_fib_fusion}), and the periodic vector of the stacked
	tensor is the sum of the periodic vectors of the two blocks at every
	positive size.
\end{theorem}

\begin{proof}
	\leanok
	\uses{thm:asymex_fib_compression, prop:asym_compression_trace}
	The word-trace identity is
	Proposition~\ref{prop:asym_compression_trace} applied to the datum of
	Theorem~\ref{thm:asymex_fib_compression}. Each matrix element of a
	periodic operator at size $L$ is the word trace at the word of length
	$L$ formed by the corresponding pairs of labels, and the periodic
	operator of the stacked tensor is the product of the two periodic
	operators, so the identity of periodic operators follows entry by entry;
	the statement for periodic vectors is the same computation.
\end{proof}

\begin{theorem}[The fusion tensors]
	\label{thm:asymex_fib_fusion_tensors}
	% Principal declaration: FibonacciCompression.fibStack_mul_fusionRight.
	\lean{FibonacciCompression.fibFusionRight, FibonacciCompression.fibFusionLeft,
		FibonacciCompression.fibonacci_right_eq, FibonacciCompression.fibonacci_left_eq,
		FibonacciCompression.fibStack_mul_fusionRight,
		FibonacciCompression.fusionLeft_mul_fibStack}
	\leanok
	\uses{def:asymex_fib, thm:asymex_fib_compression}
	The compression maps of the two slots are the fusion tensors of the
	Fibonacci algebra: $V_s$ is the column block of the adapted basis
	carrying the slot $s$, $W_s$ is the matching row block of its inverse,
	and for every letter and every slot,
	$B^{x'x}V_s=V_sC_s^{x'x}$ and $W_sB^{x'x}=C_s^{x'x}W_s$.
\end{theorem}

\begin{proof}
	\leanok
	\uses{thm:asymex_explicit_gauge}
	The identification of the two maps with the row and column blocks of
	the adapted basis is Theorem~\ref{thm:asymex_explicit_gauge}, and each
	of the two intertwining identities is an equality between products of
	matrices over $\Z[\sigma]$, transported to the complex matrices by
	Theorem~\ref{thm:asymex_golden_ring}.
\end{proof}

\begin{center}
	% Formula: B^{x'x}V_s=V_sC_s^{x'x} and W_sB^{x'x}=C_s^{x'x}W_s of
	% Theorem~\ref{thm:asymex_fib_fusion_tensors}, for $s\in\{1,\tau\}$.
	% Source: Theorem~\ref{thm:asymex_fib_fusion_tensors}.
	% Ink-to-index map: the ring is the fusion tensor $V_s$ or $W_s$;
	% the square node is the stacked source $B$ or the target block
	% $C_s\in\{B_1,B_\tau\}$.
	% Contraction set: the single virtual bond joining the ring to the
	% square node in each pair; the physical pair leg $(x',x)$ remains
	% open.
	% Boundary signature: (virtual-west=1 | virtual-east=1 |
	% physical-up=1 | physical-down=0) for each pictured equality.
	\tenkzkernel
	\begin{tenkz}[rows={wire}, cols=2, west=open, east=open]
		\tn[label pos=s, ports={90:physical:$x'x$}]{B} & \tn[skin=ring]{V_s}
	\end{tenkz}
	\quad $=$ \quad
	\begin{tenkz}[rows={wire}, cols=2, west=open, east=open]
		\tn[skin=ring]{V_s} & \tn[label pos=s, ports={90:physical:$x'x$}]{C_s}
	\end{tenkz}
	\\[2ex]
	\begin{tenkz}[rows={wire}, cols=2, west=open, east=open]
		\tn[skin=ring]{W_s} & \tn[label pos=s, ports={90:physical:$x'x$}]{B}
	\end{tenkz}
	\quad $=$ \quad
	\begin{tenkz}[rows={wire}, cols=2, west=open, east=open]
		\tn[label pos=s, ports={90:physical:$x'x$}]{C_s} & \tn[skin=ring]{W_s}
	\end{tenkz}
\end{center}

\section{Parity-graded sectors}%
\label{sec:asymex_parity}

A bond-two normal tensor paired with its own copy of weight $-1$ has the
length-$L$ coefficient $1+(-1)^L$, the power sum of the weight multiset
$\{+1,-1\}$; in the Kitaev chain this grading is the fermion-parity sign
picked up under parity-twisted boundary conditions. The point of the
example is that the source is presented in a mixing integer basis in which
neither of its two matrices has a zero row or a zero column, so the
invariant flag is not visible on the nose, and that the generic nilpotency
bound of the theorem is attained exactly.

\begin{definition}[The graded pair and the five-dimensional source]
	\label{def:asymex_parity}
	\lean{ParityGraded.parA, ParityGraded.parTargets, ParityGraded.parB}
	\leanok
	\uses{def:mps_tensor}
	The block is the bond-two tensor
	\begin{align}
		A^0 & =\begin{pmatrix}1&0\\0&0\end{pmatrix},
		    &
		A^1 & =\begin{pmatrix}0&1\\1&0\end{pmatrix},
		\notag
	\end{align}
	and the two targets are $A$ and $-A$. The source is the
	bond-dimension-five tensor
	\begin{align}
		B^0
		 & =\begin{pmatrix}
			    -1 & 1 & -1 & 1  & 3 \\
			    2  & 2 & -2 & -1 & 0 \\
			    1  & 2 & -2 & 0  & 2 \\
			    1  & 3 & -3 & -1 & 2 \\
			    -1 & 0 & 0  & 1  & 2
		    \end{pmatrix},
		 &
		B^1
		 & =\begin{pmatrix}
			    3  & 1 & 0  & -2 & -3 \\
			    -3 & 1 & -1 & 2  & 5  \\
			    -1 & 2 & -1 & 1  & 3  \\
			    0  & 1 & -1 & -1 & 1  \\
			    2  & 1 & 0  & -1 & -2
		    \end{pmatrix}.
		\notag
	\end{align}
\end{definition}

\begin{center}
	% Formula: the periodic trace tr(B^w)=(1+(-1)^{|w|})tr(A^w) of
	% Theorem~\ref{thm:asymex_parity_trace}; the source is not presented
	% as a literal stack of $A$ and $-A$, so the ring carries $B$ with
	% its coefficient, not a two-row stack.
	% Source: Definition~\ref{def:asymex_parity}.
	% Ink-to-index map: each node is the bond-dimension-five source
	% tensor $B$; the upper legs carry the letters of the word.
	% Contraction set: the horizontal virtual legs, closed around the
	% ring by the trace; no leg is left open.
	% Boundary signature: (virtual-west=0 | virtual-east=0 |
	% physical-up=3 | physical-down=0) for the finite schematic shown.
	\tenkzkernel
	\begin{tenkz}[cols=4, boundary=periodic, physical=up]
		\tn[label pos=135, ports={90:physical:$i_1$}]{B} &
		\tn[label pos=135]{B} & \tn[skin=dots]{} &
		\tn[label pos=135, ports={90:physical:$i_N$}]{B}
	\end{tenkz}
	\qquad $\tr(B^w)$, computed in
	Theorem~\ref{thm:asymex_parity_trace}
\end{center}

\begin{theorem}[Normality of the graded block]
	\label{thm:asymex_parity_normal}
	\lean{ParityGraded.parA_isNormal}
	\leanok
	\uses{def:asymex_parity, def:normal}
	The block $A$ is normal: its length-two words span $\MN{2}$.
\end{theorem}

\begin{proof}
	\leanok
	Three of the four matrix units are directly the products $A^0A^0$,
	$A^0A^1$ and $A^1A^0$, and the fourth is $A^1A^1-A^0A^0$; all four
	identities are finite verifications over integer matrices.
\end{proof}

\begin{theorem}[Compression of the graded pair]
	\label{thm:asymex_parity_compression}
	% Principal declaration: ParityGraded.parityGraded_isReduction.
	\lean{ParityGraded.parityGraded_compression,
		ParityGraded.parityGraded_isReduction,
		ParityGraded.parityGraded_left_mul_right_of_ne,
		ParityGraded.parityGraded_dim_eq,
		ParityGraded.parityGraded_evalWord_remainder_eq_zero,
		ParityGraded.parityGraded_remainder_sq_ne_zero}
	\leanok
	\uses{def:asymex_parity, thm:asym_multiblock}
	The source of Definition~\ref{def:asymex_parity} carries the block
	triangular form of Theorem~\ref{thm:asym_multiblock} for the two slots
	$A$ and $-A$ with $z=1$ zero slot and $5=2+2+1$. The compression pairs
	are mutually biorthogonal and satisfy $W_sB^wV_s=C_s^w$ for every word,
	a product of three remainder matrices vanishes, and the square of the
	remainder at the letter $0$ is nonzero, so the bound $|S|+z=3$ of
	clause (vi) is attained.
\end{theorem}

\begin{proof}
	\leanok
	\uses{thm:asymex_explicit_gauge}
	Conjugating the two letters by an explicit integer matrix with integer
	inverse produces block upper triangular matrices carrying $A$ on the
	first diagonal block, $-A$ on the second, and zero on the third, a
	finite verification over integer matrices. The compression pairs are
	then the rows and columns of that matrix and its inverse at the
	coordinates of the two slots, by
	Theorem~\ref{thm:asymex_explicit_gauge}, and computing the remainder
	from them exhibits a nonzero entry in its square.
\end{proof}

\begin{theorem}[The graded structure constant]
	\label{thm:asymex_parity_trace}
	\lean{ParityGraded.parityGraded_trace_evalWord}
	\leanok
	\uses{def:asymex_parity}
	For every nonempty word $w$,
	$\tr(B^w)=\bigl(1+(-1)^{|w|}\bigr)\tr(A^w)$.
\end{theorem}

\begin{proof}
	\leanok
	\uses{thm:asymex_parity_compression, prop:asym_compression_trace}
	Proposition~\ref{prop:asym_compression_trace} applied to the datum of
	Theorem~\ref{thm:asymex_parity_compression} gives
	$\tr(B^w)=\tr(A^w)+\tr((-A)^w)$, and pulling the weight $-1$ out of the
	second word evaluation contributes $(-1)^{|w|}\tr(A^w)$.
\end{proof}

\begin{theorem}[The weight $-1$ block has no sitewise intertwiner]
	\label{thm:asymex_parity_no_intertwiner}
	\lean{ParityGraded.parNeg_right_intertwiner_eq_zero,
		ParityGraded.parNeg_left_intertwiner_eq_zero}
	\leanok
	\uses{def:asymex_parity}
	If $V\in\C^{5\times2}$ satisfies $B^iV=V(-A^i)$ for both letters, then
	$V=0$; if $W\in\C^{2\times5}$ satisfies $WB^i=(-A^i)W$ for both
	letters, then $W=0$.
\end{theorem}

\begin{proof}
	\leanok
	Each of the two conditions is a linear system in the ten unknown
	entries. Ten of its coordinate equations, combined with explicit
	rational coefficients, determine the five entries of the first column or
	row, and each entry of the remaining one is then read off a further
	equation; both systems have only the zero solution.
\end{proof}

